\input{../tex-inputs/latex-leseansicht-vorspann}

\section[Arthur und Olga Schnitzler an Gustav Schwarzkopf, 2. 2. 1914]{L04510 Arthur und Olga Schnitzler an Gustav Schwarzkopf, 2. 2. 1914}
\nopagebreak\mylabel{L04510v}
\rehead{ }\normalsize\beginnumbering\briefempfaengerindex{Schwarzkopf, Gustav@\textsc{Schwarzkopf, Gustav}!zzzSchnitzler, Olga@\emph{von Olga Schnitzler}!1914-02-021@{2. 2. 1914}|(be}\briefempfaengerindex{Schwarzkopf, Gustav@\textsc{Schwarzkopf, Gustav}!zzzSchnitzler, Arthur@\emph{von Arthur Schnitzler}!1914-02-021@{2. 2. 1914}|(be}
\toendnotes[C]{\smallbreak\pagebreak[2]}
\correspDesc{Versand  durch Arthur Schnitzler, Olga Schnitzler am 2. 2. 1914 in Salzburg
\newline{}Übermittlung  am 2. 2. 1914 in Salzburg
\newline{}Erhalt  durch Gustav Schwarzkopf im Zeitraum [3. 2. 1914 – 7. 2. 1914?] in Wien}\toendnotes[C]{\smallbreak}
\Standort{DLA, A:Schnitzler, HS.1985.1.1897.}
\physDesc{Bildpostkarte, 152 Zeichen
\newline{}Handschrift Arthur Schnitzler: Bleistift, deutsche Kurrent
\newline{}Handschrift Olga Schnitzler: Bleistift, lateinische Kurrent
\newline{}Versand: Stempel: »\nobreak{}\oindex{Salzburg@\textbf{Salzburg}|pwk}Salzburg 2, 2. II. 14, 9\nobreak{}«.  }
\pstart
           \noindent{}\raggedleft{}{\pb}\textcolor{gray}{\textbf{Salzburg. Makartplatz\oindex{Makartplatz@\textbf{Makartplatz}|pw} mit Dreifaltigkeitskirche\oindex{Church of the Holy Trinity, Salzburg@\textbf{Church of the Holy Trinity, Salzburg}|pw}}}\pend
           \pstart{}{\pb}Hrn. \textsc{Gustav Schwarzkopf}\pend{}\pstart{}Wien I\oindex{I., Innere Stadt@\textbf{I., Innere Stadt}|pw}\pend{}\pstart{}\textsc{Tiefer Graben 17}\oindex{Wien@\textbf{Wien}!I., Innere Stadt@\textbf{I., Innere Stadt}!Tiefer Graben 17@\textbf{Tiefer Graben 17}|pw}\pend{}{\bigskip}\vspace{1em}
\pstart
           \raggedleft{}2. 2. 14\pend
           \vspace{0.5em}
\pstart
           Herzliche Grüße! Dienſtag Abend hoffen wir daheim
      zu ſein. Auf baldiges
      Wiederſehn!\pend
           \pstart Ihr \spacefill\mbox{A.}\pend{}\selectlanguage{ngerman}\vspace{1em}
\pstart
           \noindent{}{[}hs. O. Schnitzler:{]} Herzlichst \spacefill\mbox{Olga}\pend
           \selectlanguage{ngerman}\endnumbering\briefempfaengerindex{Schwarzkopf, Gustav@\textsc{Schwarzkopf, Gustav}!zzzSchnitzler, Olga@\emph{von Olga Schnitzler}!1914-02-021@{2. 2. 1914}|)be}\briefempfaengerindex{Schwarzkopf, Gustav@\textsc{Schwarzkopf, Gustav}!zzzSchnitzler, Arthur@\emph{von Arthur Schnitzler}!1914-02-021@{2. 2. 1914}|)be}\mylabel{L04510h}
\begin{anhang}
\end{anhang}\newcommand{\dateiname}{L04510}\newcommand{\titel}{Arthur und Olga Schnitzler an Gustav Schwarzkopf, 2. 2. 1914}\newcommand{\editorInnen}{Herausgegeben von Jahnke, SelmaMüller, Martin Anton}\input{../tex-inputs/latex-leseansicht-abspann}
